% Principe de la tomographie par émission de positons (exercice 12, chapitre 6)
\begin{tikzpicture}[x=1mm,y=1mm]
\foreach \a in {-33,-23,...,317}{%
  \filldraw[fill=otNavyTint,draw=otNavy!60,line width=0.4pt,rotate=\a] (-2,24) rectangle (2,29);}
\filldraw[fill=otRed!35,draw=otRed!80!black,line width=0.5pt,rotate=-33] (-2,24) rectangle (2,29);
\filldraw[fill=otRed!35,draw=otRed!80!black,line width=0.5pt,rotate=157] (-2,24) rectangle (2,29);
\filldraw[fill=otOrange!12,draw=otOrange!70!black,line width=0.6pt] (0,0) ellipse (13 and 16);
\node[pcLegende,text=otOrange!70!black] at (-4.5,9) {patient};
\draw[otRed!80!black,line width=1pt,-{Stealth[length=2.2mm]}] (4,3) -- ++(62:21);
\draw[otRed!80!black,line width=1pt,-{Stealth[length=2.2mm]}] (4,3) -- ++(242:29.5);
\fill[otGreenDark] (4,3) circle (1.1);
\pcNote[west]{(5,3.5)}{(40,8)}{annihilation d'un positon et d'un électron}
\pcNote[west]{(11.9,17.9)}{(40,20)}{photon $\gamma$ de 511 keV}
\pcNote[west]{(28.5,0)}{(40,-4)}{anneau de détecteurs}
\pcNote[west]{(-7.3,-18.2)}{(40,-16)}{détections simultanées : l'annihilation a eu lieu sur ce segment}
\end{tikzpicture}
\par\smallskip\legende{Les deux photons $\gamma$, émis en sens opposés, sont détectés au même instant par deux détecteurs de l'anneau.}
